\section{推理时搜索：从单轨执行到可回溯决策}

\subsection{为什么必须引入搜索}
讲者给出的结论非常明确：没有搜索能力，Agent 一旦走偏就很难回到正确轨道。尤其在 Web 场景中，错误动作会改变页面状态，后续操作空间随之变化，恢复成本极高。

\begin{importantbox}{搜索的直接收益}
\begin{itemize}
    \item 允许多候选轨迹并行评估；
    \item 支持回溯与分支重试；
    \item 把 ``一次猜对'' 变成 ``多次试探后选优''。
\end{itemize}
\end{importantbox}

\subsection{基础策略：Rejection Sampling}
最简单的做法是先生成多条轨迹，再用终局标准筛选。它实现简单、并行友好，但效率不高，因为大量无效轨迹会浪费推理预算。

\begin{knowledgebox}{Rejection Sampling 何时有效}
当任务较短、验证器可靠、并行资源充足时，这种方法可以作为低工程成本基线，并快速带来增益。
\end{knowledgebox}

\subsection{进阶策略：价值引导与树搜索}
当任务长度增加，纯采样很快失效。此时要引入价值函数或 Judge 估计中间状态潜力，优先扩展高价值节点，实现 ``边探索边裁枝''。

\begin{figure}[H]
\centering
\includegraphics[width=0.82\textwidth]{frame-07-search-eval.jpg}
\caption{视频约 00:34:39：讲者说明在搜索分支中评估轨迹结果并做继续/终止决策}
\end{figure}

\begin{table}[H]
\centering
\small
\begin{tabular}{p{2.8cm}p{3.8cm}p{4.4cm}}
\toprule
\textbf{方法} & \textbf{优点} & \textbf{不足} \\
\midrule
单轨执行 & 延迟低、实现简单 & 错误不可恢复，成功率受单次决策噪声影响大 \\
Rejection Sampling & 并行友好，容易上线 & 计算浪费明显，长任务收益下降 \\
价值引导搜索 & 样本效率高，可回溯 & 工程复杂，对 Judge 质量敏感 \\
\bottomrule
\end{tabular}
\caption{Agent 推理策略对比}
\end{table}

\begin{warningbox}{搜索不是免费午餐}
搜索会显著增加推理成本。若没有良好剪枝和节点复用机制，系统可能陷入 ``预算耗尽但仍未收敛'' 的状态。
\end{warningbox}

\subsection{本章小结}
在长链路网页任务中，搜索能力不是可选项，而是鲁棒性的核心来源。关键不在于 ``是否搜索''，而在于 ``如何高效搜索''。

